%%%%%%%%%%%%%%%%%%%%%%%%%%%%%%%%%%%%%%%%%
% Zhovten Games Game Studio Article
% LaTeX Template
% Version 1.7 (April 19, 2026)
%
% This template was created by:
% Vel (enquiries@latextypesetting.com)
% LaTeXTypesetting.com
%
% Modified by Sam Starling (www.linkedin.com/in/sam-starling/, ORCID: https://orcid.org/0009-0009-2621-6372)
%
%!TEX program = xelatex
% Note: this template must be compiled with XeLaTeX rather than PDFLaTeX
% due to the custom fonts used. The line above should ensure this happens
% automatically, but if it doesn't, your LaTeX editor should have a simple toggle
% to switch to using XeLaTeX.
%
%%%%%%%%%%%%%%%%%%%%%%%%%%%%%%%%%%%%%%%%%

%----------------------------------------------------------------------------------------
%	CLASS, PACKAGES AND OTHER DOCUMENT CONFIGURATIONS
%----------------------------------------------------------------------------------------

% Options for packages loaded elsewhere
\PassOptionsToPackage{unicode}{hyperref}
% \PassOptionsToPackage{hyphens}{xurl}
% 
\documentclass{ZGJournalArticle_2026}

\defaultfontfeatures{Scale=MatchLowercase}

\IfFileExists{upquote.sty}{\usepackage{upquote}}{}
\IfFileExists{microtype.sty}{% use microtype if available
  \usepackage[]{microtype}
  \UseMicrotypeSet[protrusion]{basicmath} % disable protrusion for tt fonts
}{}

% A wrapper around soul necessary to make underlining work for Greek.
\usepackage{xesoul}

\usepackage{xcolor}
\usepackage{calc} % for calculating minipage widths
\usepackage[obeyspaces,spaces,hyphens]{xurl}
\urlstyle{same} % Disable monospaced font for ordinary URLs.

% xurl already provides broad line-breaking rules for URLs and paths.
% Do not override \UrlBreaks manually: long path segments may require
% breakpoints beyond '/', '-', '.', and '_'.

\usepackage{fvextra} % Controlled wrapping for verbatim code blocks.

\renewcommand{\floatpagefraction}{.8}%
\setlength{\emergencystretch}{3em} % prevent overfull lines


%----------------------------------------------------------------------------------------
%	PANDOC COMPATIBILITY
%----------------------------------------------------------------------------------------

% Pandoc lists
\providecommand{\tightlist}{%
  \setlength{\itemsep}{0pt}\setlength{\parskip}{0pt}}

% Pandoc tables
% Correct order of tables after \paragraph or \subparagraph.
\usepackage{etoolbox}
\makeatletter
\patchcmd\longtable{\par}{\if@noskipsec\mbox{}\fi\par}{}{}
\makeatother

% Pandoc may emit \def\LTcaptype{none} for captionless longtable blocks.
% Some caption/longtable combinations still look for a counter named `none`.
% Define a dummy counter so uncaptioned Pandoc longtables do not fail.
\makeatletter
\@ifundefined{c@none}{\newcounter{none}}{}
\makeatother

% Allow footnotes in longtable head/foot.
\IfFileExists{footnotehyper.sty}{\usepackage{footnotehyper}}{\usepackage{footnote}}
\makesavenoteenv{longtable}

% Pandoc task lists
% Pandoc may emit \boxtimes / \square for GitHub-style task list checkboxes.
\usepackage{amssymb}

% Pandoc images
\usepackage{graphicx}
\makeatletter
\newsavebox\pandoc@box
\newcommand*\pandocbounded[1]{% scales image to fit in text height/width
  \sbox\pandoc@box{#1}%
  \Gscale@div\@tempa{\textheight}{\dimexpr\ht\pandoc@box+\dp\pandoc@box\relax}%
  \Gscale@div\@tempb{\linewidth}{\wd\pandoc@box}%
  \ifdim\@tempb\p@<\@tempa\p@\let\@tempa\@tempb\fi% select the smaller of both
  \ifdim\@tempa\p@<\p@\scalebox{\@tempa}{\usebox\pandoc@box}%
  \else\usebox{\pandoc@box}%
  \fi%
}
\def\maxwidth{\ifdim\Gin@nat@width>\linewidth\linewidth\else\Gin@nat@width\fi}
\def\maxheight{\ifdim\Gin@nat@height>\textheight\textheight\else\Gin@nat@height\fi}
\makeatother

% Scale images if necessary, so that they will not overflow the page
% margins by default, and it is still possible to overwrite the defaults
% using explicit options in \includegraphics[width, height, ...]{}.
\setkeys{Gin}{width=\maxwidth,height=\maxheight,keepaspectratio}

% Set default figure placement to htbp.
\makeatletter
\def\fps@figure{htbp}
\makeatother

% Pandoc citations
% CSL/citeproc support is kept unconditional here because the build pipeline
% can inject pre-rendered bibliography sections that already contain
% \begin{CSLReferences} ... \end{CSLReferences} even when the final pandoc pass
% itself is executed without --citeproc (for example with suppress-bibliography=true).
% Without these definitions XeLaTeX fails with:
%   LaTeX Error: Environment CSLReferences undefined.
% In other words, the bibliography block is already in the payload, so the
% template must always know how to typeset it.
\NewDocumentCommand\citeproctext{}{}
\NewDocumentCommand\citeproc{mm}{%
  \begingroup\def\citeproctext{#2}\cite{#1}\endgroup}
\makeatletter
% Allow citations to break across lines.
\let\@cite@ofmt\@firstofone
% Avoid brackets around text for \cite.
\def\@biblabel#1{}
\def\@cite#1#2{{#1\if@tempswa , #2\fi}}
\makeatother
\newlength{\cslhangindent}
\setlength{\cslhangindent}{1.5em}
\newlength{\csllabelwidth}
\setlength{\csllabelwidth}{3em}
\newenvironment{CSLReferences}[2] % #1 hanging-indent, #2 entry-spacing
 {\begin{list}{}{% This helps avoid over-stretched lines for DOI links in bibliography.
  \sloppy
  \setlength{\itemindent}{0pt}
  \setlength{\leftmargin}{0pt}
  \setlength{\parsep}{0pt}
  % Turn on hanging indent if param 1 is 1.
  \ifodd #1
   \setlength{\leftmargin}{\cslhangindent}
   \setlength{\itemindent}{-1\cslhangindent}
  \fi
  % Set entry spacing.
  \setlength{\itemsep}{#2\baselineskip}}}
 {\end{list}}
\newcommand{\CSLBlock}[1]{\linebreak \hfill\break\parbox[t]{\linewidth}{\strut\ignorespaces#1\strut}}
\newcommand{\CSLLeftMargin}[1]{\parbox[t]{\csllabelwidth}{\strut#1\strut}}
\newcommand{\CSLRightInline}[1]{\parbox[t]{\linewidth - \csllabelwidth}{\strut#1\strut}}
\newcommand{\CSLIndent}[1]{\hspace{\cslhangindent}#1}

% Pandoc code highlighting
% Required when fenced code blocks are converted to LaTeX with highlighting enabled.
% Defines environments/macros such as Shaded, Highlighting, NormalTok, KeywordTok, etc.
\usepackage{color}
\usepackage{fancyvrb}
\newcommand{\VerbBar}{|}
\newcommand{\VERB}{\Verb[commandchars=\\\{\}]}
\DefineVerbatimEnvironment{Highlighting}{Verbatim}{commandchars=\\\{\}}
% Add ',fontsize=\small' for more characters per line
\newenvironment{Shaded}{}{}
\newcommand{\AlertTok}[1]{\textcolor[rgb]{1.00,0.00,0.00}{\textbf{#1}}}
\newcommand{\AnnotationTok}[1]{\textcolor[rgb]{0.38,0.63,0.69}{\textbf{\textit{#1}}}}
\newcommand{\AttributeTok}[1]{\textcolor[rgb]{0.49,0.56,0.16}{#1}}
\newcommand{\BaseNTok}[1]{\textcolor[rgb]{0.25,0.63,0.44}{#1}}
\newcommand{\BuiltInTok}[1]{\textcolor[rgb]{0.00,0.50,0.00}{#1}}
\newcommand{\CharTok}[1]{\textcolor[rgb]{0.25,0.44,0.63}{#1}}
\newcommand{\CommentTok}[1]{\textcolor[rgb]{0.38,0.63,0.69}{\textit{#1}}}
\newcommand{\CommentVarTok}[1]{\textcolor[rgb]{0.38,0.63,0.69}{\textbf{\textit{#1}}}}
\newcommand{\ConstantTok}[1]{\textcolor[rgb]{0.53,0.00,0.00}{#1}}
\newcommand{\ControlFlowTok}[1]{\textcolor[rgb]{0.00,0.44,0.13}{\textbf{#1}}}
\newcommand{\DataTypeTok}[1]{\textcolor[rgb]{0.56,0.13,0.00}{#1}}
\newcommand{\DecValTok}[1]{\textcolor[rgb]{0.25,0.63,0.44}{#1}}
\newcommand{\DocumentationTok}[1]{\textcolor[rgb]{0.73,0.13,0.13}{\textit{#1}}}
\newcommand{\ErrorTok}[1]{\textcolor[rgb]{1.00,0.00,0.00}{\textbf{#1}}}
\newcommand{\ExtensionTok}[1]{#1}
\newcommand{\FloatTok}[1]{\textcolor[rgb]{0.25,0.63,0.44}{#1}}
\newcommand{\FunctionTok}[1]{\textcolor[rgb]{0.02,0.16,0.49}{#1}}
\newcommand{\ImportTok}[1]{\textcolor[rgb]{0.00,0.50,0.00}{\textbf{#1}}}
\newcommand{\InformationTok}[1]{\textcolor[rgb]{0.38,0.63,0.69}{\textbf{\textit{#1}}}}
\newcommand{\KeywordTok}[1]{\textcolor[rgb]{0.00,0.44,0.13}{\textbf{#1}}}
\newcommand{\NormalTok}[1]{#1}
\newcommand{\OperatorTok}[1]{\textcolor[rgb]{0.40,0.40,0.40}{#1}}
\newcommand{\OtherTok}[1]{\textcolor[rgb]{0.00,0.44,0.13}{#1}}
\newcommand{\PreprocessorTok}[1]{\textcolor[rgb]{0.74,0.48,0.00}{#1}}
\newcommand{\RegionMarkerTok}[1]{#1}
\newcommand{\SpecialCharTok}[1]{\textcolor[rgb]{0.25,0.44,0.63}{#1}}
\newcommand{\SpecialStringTok}[1]{\textcolor[rgb]{0.73,0.40,0.53}{#1}}
\newcommand{\StringTok}[1]{\textcolor[rgb]{0.25,0.44,0.63}{#1}}
\newcommand{\VariableTok}[1]{\textcolor[rgb]{0.10,0.09,0.49}{#1}}
\newcommand{\VerbatimStringTok}[1]{\textcolor[rgb]{0.25,0.44,0.63}{#1}}
\newcommand{\WarningTok}[1]{\textcolor[rgb]{0.38,0.63,0.69}{\textbf{\textit{#1}}}}

% Allow fenced code blocks to wrap long lines without adding visible
% continuation markers. The PDF remains visually clean and code remains
% selectable as text.
\fvset{
  breaklines=true,
  breakanywhere=true,
  breaksymbolleft={},
  breaksymbolright={}
}

% Apply brand-level styles after Pandoc has defined the Shaded environment.
\zgapplycontentstyles

\makeatletter
\@ifpackageloaded{subcaption}{}{\usepackage{subcaption}}
\@ifpackageloaded{caption}{}{\usepackage{caption}}
\captionsetup[subfigure]{margin=0.5em}
\AtBeginDocument{%
\renewcommand*\figurename{Figure}
\renewcommand*\tablename{Table}
}
\AtBeginDocument{%
\renewcommand*\listfigurename{List of Figures}
\renewcommand*\listtablename{List of Tables}
}
\newcounter{pandoccrossref@subfigures@footnote@counter}
\newenvironment{pandoccrossrefsubfigures}{%
\setcounter{pandoccrossref@subfigures@footnote@counter}{0}
\begin{figure}\centering%
\gdef\global@pandoccrossref@subfigures@footnotes{}%
\DeclareRobustCommand{\footnote}[1]{\footnotemark%
\stepcounter{pandoccrossref@subfigures@footnote@counter}%
\ifx\global@pandoccrossref@subfigures@footnotes\empty%
\gdef\global@pandoccrossref@subfigures@footnotes{{##1}}%
\else%
\g@addto@macro\global@pandoccrossref@subfigures@footnotes{, {##1}}%
\fi}}%
{\end{figure}%
\addtocounter{footnote}{-\value{pandoccrossref@subfigures@footnote@counter}}
\@for\f:=\global@pandoccrossref@subfigures@footnotes\do{\stepcounter{footnote}\footnotetext{\f}}%
\gdef\global@pandoccrossref@subfigures@footnotes{}}
\@ifpackageloaded{float}{}{\usepackage{float}}
\floatstyle{ruled}
\@ifundefined{c@chapter}{\newfloat{codelisting}{h}{lop}}{\newfloat{codelisting}{h}{lop}[chapter]}
\floatname{codelisting}{Listing}
\newcommand*\listoflistings{\listof{codelisting}{List of Listings}}
\makeatother

%----------------------------------------------------------------------------------------
%	ARTICLE INFORMATION
%----------------------------------------------------------------------------------------

% Set the page number

\usepackage{polyglossia}
\setdefaultlanguage{en-GB}



\setzgbrandassetsdir{C:/Users/SamTykhon/Documents/Persons/Sam/Projects/prod/Gamedev/Zhovten
Games/zg-journal-template/template/branding/default}

\setzglogo{LogoZG.png}

\setzgjournalname{Game Studio Zhovten Games}



\setzgjournalurl{https://zhovten-games.github.io/}


\setzgprimarycolorhtml{F28C28}

\setzgabstractbggray{0.93}

\setzgcodebggray{0.965}

\setzgquotebggray{0.985}

\setzgquoteframegray{0.72}

\setzgheaderlefttext{Zhovten Games Journal}

\hypersetup{
		pdftitle={Literate Programming: Donald Knuth{,} WEB{,} and Contemporary Workflows},
      	pdflang={en},
      	colorlinks=true, % Whether to color the text of links
	urlcolor=ZGprimary, % Color for \url and \href links
	linkcolor=ZGprimary, % Color for \nameref links
	citecolor=ZGprimary, % Color of reference citations
  pdfcreator={LaTeX via pandoc}}

\zgtitle{Literate Programming: Donald Knuth, WEB, and Contemporary
Workflows} % Article title, to appear on the title page and in page headers

\zgheadertitle{Literate Programming: Donald Knuth, WEB, and Contemporary
Workflows} % If the shorter title is not available, use the default one

\zgccbox{\includegraphics[width=2.25cm]{cc-by-sa.pdf}}

\zgorcidlogo{
  \includegraphics[width=0.3cm]{ORCIDiD.pdf}
}

\zgauthors{ % Specify all article authors within this command
	% Each article author needs to be in a \articleauthor command, the first argument of which is the author's name and the second the institution and contact information
      \articleauthor{Sam Starling}{Zhovten
Games, \href{mailto:sam@phantom-draft.com}{sam@phantom-draft.com}}{\zgorcidlogo\ \url{0009-0009-2621-6372}}
    \articleauthor{Oksana Dubinetska}{Zhovten
Games, \href{mailto:oksana.dubinetska@phantom-draft.com}{oksana.dubinetska@phantom-draft.com}}{\zgorcidlogo\ \url{0009-0003-8777-8412}}
    
}

\zgauthorsheaders{} % Legacy single-author fallback.

\zgvolumeandissue{}
\zgvolume{}

\zgyear{}

\zgdoi{}

\zgabstract{Donald Knuth's 1984 article \emph{Literate Programming}
proposed that programs should be written primarily for human readers and
only secondarily transformed into machine-executable form. This review
reconstructs that argument through the WEB system and examines how its
central principle persists across contemporary engineering practices. It
distinguishes direct or historically proximate continuations, such as
CWEB and noweb, from instrumentally related workflows, such as Org
Babel, and from broader conceptual associations involving executable
documents, reproducible environments, IDE navigation, and LLM-assisted
development. A companion suite remasters Knuth's prime-number example in
several deterministic workflows and introduces Prompt-Literate Workflow
as a controlled boundary case for probabilistic generation. The review
argues that LLM output should be treated not as a build artifact, but as
a candidate artifact governed by a human-authored plan, chunk contracts,
bounded prompts, review, tests, smoke-checks, and traceability. The
concluding section extends this logic towards literate design and GDD
pipelines.}

\zgkeywords{} % Keywords

%----------------------------------------------------------------------------------------

\begin{document}

\titlesection % Output the article title section, automatically populated using the information in the ARTICLE INFORMATION above

%----------------------------------------------------------------------------------------
%	ARTICLE BODY
%----------------------------------------------------------------------------------------
\subsection{Methodological Framework of the
Review}\label{methodological-framework-of-the-review}

The breadth of the issues addressed is not accidental. Donald Knuth's
article is not narrowly focused documentation for the WEB system, but a
methodological text written by a practising programmer who speaks
simultaneously as an engineer and as a theorist of the profession. In
discussing a specific tool, Knuth addresses the structure of a program,
the order in which it should be explained, style, portability,
development costs, the publication of code, and the future form of
programming environments. Any contemporary commentary on this work
therefore inevitably extends beyond the reconstruction of WEB and
requires several engineering lineages to be compared.

At the same time, the tools and practices discussed in this review are
not treated as a single genealogical tree in which every modern
mechanism literally descends from WEB. For the purposes of the
subsequent analysis, it is important to distinguish three types of
relationship. The first is a direct or historically proximate
continuation of WEB: above all, CWEB and noweb. The second is a practice
with an analogous instrumental logic: for example, Org Babel tangling,
which employs weaving/tangling mechanisms, although this alone does not
establish a direct genealogical connection. The third is a conceptual
association, or a productive comparison: WEB ↔ notebooks, WEB ↔
workflows based on large language models (LLMs)\footnote{A large
  language model, or LLM, is a statistical model trained to continue and
  transform text according to context and instructions.}, WEB ↔
integrated development environment (IDE)\footnote{An integrated
  development environment, or IDE, combines editing, navigation,
  execution, debugging, and other development tools.} navigation and
context maps. In the latter case, the issue is the recurrence of a
similar problem: how to connect code, explanation, reading order,
navigation, and a verifiable result.

This comparison follows the logic that Donald Knuth himself sets out in
the Related Work section: the strength of WEB lay not in the invention
of each individual element, but in bringing existing ideas together into
an integrated working method.

\subsection{Table of Contents}\label{table-of-contents}

\subsubsection{Main Article}\label{main-article}

\begin{itemize}
\tightlist
\item
  \textbf{Opening Remarks} --- the historical framing of the problem,
  the comparative framework, and the rationale for a contemporary
  remaster approach.
\item
  \textbf{A. Introduction} --- a shift in whom programming is primarily
  addressed to: from instructing the machine to explaining the program
  to a human reader.
\item
  \textbf{B. The WEB System} --- WEB as a system that generates
  human-oriented and machine-oriented projections from a single source.
\item
  \textbf{C--F. Example / source / tangle / weave} --- the prime-number
  example, the original WEB file, the machine-oriented TANGLE
  projection, and the human-oriented WEAVE projection.
\item
  \textbf{G. Additional Bells and Whistles} --- WEB's additional
  capabilities and the transition from an instructional example to
  production-level complexity.
\item
  \textbf{H. Occam's Razor} --- the limits of the tool and the need for
  methodological restraint.
\item
  \textbf{I. Portability} --- the portability of the literate approach
  beyond the original TeX / Pascal pairing.
\item
  \textbf{J. Programs as Webs} --- the program as a network of named
  fragments, dependencies, and explanatory relationships.
\item
  \textbf{K. Stylistic Issues} --- style, section naming, and the
  readability of literate source.
\item
  \textbf{L. Economic Issues} --- the costs of introducing, maintaining,
  and consistently applying literate programming.
\item
  \textbf{M. Related Work} --- related lineages: structured programming,
  documentation systems, and executable-document practices.
\item
  \textbf{N. Retrospect and Prospects} --- a retrospective assessment of
  the method and its future prospects.
\item
  \textbf{Conclusions 1--5} --- which aspects of Knuth's idea persist
  literally, which have shifted towards reproducible research and
  notebook environments, and which are re-emerging in LLM-assisted
  workflows.
\end{itemize}

\subsubsection{Repository Navigation}\label{repository-navigation}

\paragraph{Canonical Source of the
Article}\label{canonical-source-of-the-article}

\begin{itemize}
\tightlist
\item
  public repository root:
  {\urlstyle{tt}\path{https://github.com/Zhovten-Games/literate-programming}};
\item
  canonical source of the review:
  {\urlstyle{tt}\path{01-Literate Programming - Knuth WEB and Contemporary Workflows/01-Literate Programming - Knuth WEB and Contemporary Workflows.md}};
\item
  author metadata file:
  {\urlstyle{tt}\path{01-Literate Programming - Knuth WEB and Contemporary Workflows/authors.yml}};
\item
  bibliography inventory layer:
  {\urlstyle{tt}\path{01-Literate Programming - Knuth WEB and Contemporary Workflows/Bibliography/Original.md}};
\item
  sectioned bibliography JSON:
  {\urlstyle{tt}\path{01-Literate Programming - Knuth WEB and Contemporary Workflows/Bibliography/*.json}}.
\end{itemize}

\paragraph{Companion Suite}\label{companion-suite}

\begin{itemize}
\tightlist
\item
  root of the companion suite:
  {\urlstyle{tt}\path{02-Literate-Companion-Implementations/}};
\item
  shared README for the companion suite:
  {\urlstyle{tt}\path{02-Literate-Companion-Implementations/README.md}};
\item
  Makefile for verification runs:
  {\urlstyle{tt}\path{02-Literate-Companion-Implementations/Makefile}};
\item
  English branch of the examples:
  {\urlstyle{tt}\path{02-Literate-Companion-Implementations/examples/en/}};
\item
  Russian branch of the examples:
  {\urlstyle{tt}\path{02-Literate-Companion-Implementations/examples/ru/}}.
\end{itemize}

\paragraph{Deterministic Companion Examples
01--06}\label{deterministic-companion-examples-0106}

\begin{itemize}
\tightlist
\item
  \texttt{01-cweb} --- the more canonical CWEB route, with an explicit
  two-branch {\urlstyle{tt}\path{tangle / weave}} model;
\item
  \texttt{02-noweb-like} --- the primary lightweight C++ remaster,
  centred on a \texttt{tangle-first\ workflow};
\item
  \texttt{03-org-babel} --- a plain-text literate workflow in an
  Org-mode / Emacs environment, with tangling;
\item
  \texttt{04-quarto} --- an executable-document route for
  publication-ready HTML / computational publishing;
\item
  \texttt{05-jupyter} --- a notebook route with particular attention to
  execution, hidden state, and reproducible execution;
\item
  \texttt{06-rmarkdown} --- a reproducible report workflow based on R
  Markdown / knitr / Pandoc.
\end{itemize}

\paragraph{Prompt-Literate Workflow and the Boundary LLM
Example}\label{prompt-literate-workflow-and-the-boundary-llm-example}

\begin{itemize}
\tightlist
\item
  public repository for the Prompt-Literate Workflow methodology:
  {\urlstyle{tt}\path{https://github.com/IRONCREED/prompt-literate-workflow}};
\item
  role: a reusable methodology repository and GitHub template
  repository;
\item
  integration in \texttt{literate-programming}: a root
  submodule/scaffold mounted at
  {\urlstyle{tt}\path{prompt-literate-workflow/}};
\item
  methodology path used by the companion examples:
  {\urlstyle{tt}\path{02-Literate-Companion-Implementations/methodology/prompt-literate-workflow/}};
\item
  article-specific local specialisation for Knuth's prime-number
  problem:
  {\urlstyle{tt}\path{02-Literate-Companion-Implementations/methodology/extensions/primes-example/}};
\item
  \texttt{07-prompt-literate} --- a demonstrative application of
  Prompt-Literate Workflow (PLW) to the prime-number example, rather
  than a definition of the method itself.
\end{itemize}

\paragraph{\texorpdfstring{Working Files within
\texttt{07-prompt-literate}}{Working Files within 07-prompt-literate}}\label{working-files-within-07-prompt-literate}

\begin{itemize}
\tightlist
\item
  \texttt{COMPANION.md} --- entry point for the companion example;
\item
  \texttt{primes.plan.md} --- human-authored plan and canonical source
  for the prompt-literate example;
\item
  \texttt{CONTRACTS.md} --- chunk contracts for the prime-number task;
\item
  \texttt{SCENARIOS.md} --- validation scenarios and acceptance
  criteria;
\item
  {\urlstyle{tt}\path{prompts/fill-chunks.prompt.md}} --- a bounded
  prompt for filling only the permitted \texttt{LLM-TODO} chunks;
\item
  {\urlstyle{tt}\path{prompts/review-generated-code.prompt.md}} --- a
  prompt for review, rather than for rewriting the architecture;
\item
  {\urlstyle{tt}\path{generated/.gitkeep}} --- an empty directory for a
  future accepted generated artifact;
\item
  {\urlstyle{tt}\path{tests/smoke-check.sh}} --- verification of the
  accepted generated code;
\item
  \texttt{output.expected.txt} --- expected output markers;
\item
  \texttt{TRACE.md} --- a log of the model, prompt run, review, manual
  edits, validation, and acceptance decision.
\end{itemize}

\section{Main Article}\label{main-article-1}

\subsection{Opening Remarks}\label{opening-remarks}

The historical significance of Donald Knuth's article \emph{Literate
Programming} becomes particularly clear in retrospect. The text was
published in 1984, before the later canonisation of pattern language in
\emph{Design Patterns} (1994)
(\citeproc{ref-gamma1994designPatterns}{Gamma et al. 1994}) and long
before \emph{Clean Architecture} (2017)
(\citeproc{ref-martin2017cleanArchitecture}{Martin 2017}). In this
sense, it represents one of the early attempts to impose methodological
order on the complexity of software development not only through the
structure of code, but also through the structure of its explanation.

\emph{Literate Programming} can be read as a precise formulation of a
problem that remains relevant: a program must be not only executable but
also explainable. Knuth proposes changing whom programming is primarily
addressed to: rather than treating the instruction of a computer as the
primary task, a programmer should first explain to a human reader what
the computer is intended to do (\citeproc{ref-knuth1984a}{Knuth 1984}).
This shift does not negate machine execution, but subordinates it to the
broader task of organising a program in an order suited to human
understanding.

In 2026, this formulation has acquired additional relevance because the
software-development workflow itself has changed. If the classical model
can be described as a movement from one human being to another and then
to a machine, an AI-assisted workflow increasingly involves a chain in
which a human formulates an intention for a machine, while one
computational system helps produce input for another. A prompt, however,
is not in itself a literate source. A contemporary AI-based approach can
be regarded as a substantive continuation of literate programming only
if it preserves an engineering discipline: explanation → specification →
code → tests → artifact. Without this chain, what accelerates is not so
much the understanding of a program as the production of code with
insufficient verifiability and a weak explanatory structure.

In this review, WEB is treated as the original form of this idea. In
Donald Knuth's model, a single source generates two projections: a
document for the human reader and a program for the machine. This
relationship between explanation, code, and the build mechanism
connecting them serves as the principal criterion for the analysis that
follows. The purpose of the review is not to reproduce the Pascal
example as a self-contained historical object. The example's codebase
has been remastered in C++, while the primary demonstration route has
been implemented in a noweb-like form: it preserves named chunks and
tangling without requiring entry into the more elaborate CWEB/TeX
toolchain.

The subsequent analysis compares several instrumental and methodological
lineages:

\begin{itemize}
\tightlist
\item
  \textbf{CWEB} (\citeproc{ref-knuth1992a}{Knuth and Levy, n.d.}),
\item
  \textbf{noweb-like C++} (\citeproc{ref-ramsey1994a}{Ramsey, n.d.}),
\item
  \textbf{Org Babel}
  (\citeproc{ref-orgmanualExtractingSourceCode}{{``Org Manual:
  Extracting Source Code,''} n.d.}),
\item
  \textbf{Quarto / Jupyter / R Markdown}
  (\citeproc{ref-quartoDocs}{{``Quarto Documentation,''} n.d.};
  \citeproc{ref-jupyterDocs}{{``Project Jupyter Documentation,''} n.d.};
  \citeproc{ref-rmarkdownDocs}{{``R Markdown Documentation,''} n.d.}).
\end{itemize}

The technical branches are documented in companion materials within the
code repository. The main text follows Donald Knuth's article and tests
its central proposition: whether a program can first be organised in an
order suited to human understanding and only then transformed into
machine-executable code.

The conclusion returns to several questions: how literate programming
relates to the later language of engineering patterns; which of Knuth's
claims were realised literally; which ideas shifted into the domain of
reproducible research and notebook environments; and what is now
re-emerging around LLMs. Particular attention will be given to the
proposition that an LLM should receive not an unconstrained request, but
a bounded operation over a human-authored literate plan, chunk
contracts, bounded prompts, tests/smoke-check, and TRACE.

There is no full-fledged \texttt{07-llm} companion here that would be
equivalent in reproducibility to CWEB, noweb, Org Babel, Quarto,
Jupyter, or R Markdown: LLM generation depends on the model, version,
context, and execution mode. Instead, a boundary example,
\texttt{07-prompt-literate}, has been added to demonstrate the
application of Prompt-Literate Workflow. Its status and limitations will
be examined in detail in the conclusion.

Prompt-Literate Workflow was initially formulated within the present
review and subsequently refined through its practical application to an
independent core component of a game project. This pilot demonstrated
the need to separate the method's compact, universal core from a
project-local extension layer. Domain-specific rules must not silently
override the basic invariants: they should be formalised as local
extensions, while the inclusion of a local rule in the core requires a
separate process of generalisation and promotion into the methodology's
foundational layer.

\subsection{A. Introduction}\label{a.-introduction}

Donald Knuth begins by establishing a professional norm. Structured
programming has already made programs more reliable and easier to
understand, but this is not sufficient. He sees the next step not in the
introduction of a new control construct, but in changing the way a
program is presented.

The image of the programmer as an essayist is not merely decorative.
Understanding must be embedded in the structure of the program itself:
concepts are introduced in an order suited to the reader, while formal
and informal means of exposition reinforce one another.

\subsection{B. The WEB System}\label{b.-the-web-system}

In the second section, Donald Knuth demonstrates that \emph{literate
programming} is not a metaphor, but a concrete system: WEB. Its purpose
is to combine a document language with a programming language so that a
single source can generate two distinct results: a readable description
for a human reader and an executable program for a machine.

Donald Knuth proposes that a complex program should be understood as a
network of simple parts and the relationships between them.

\textbf{It is useful here to distinguish three terms:}

\begin{itemize}
\tightlist
\item
  \textbf{WEB} --- from the English word \emph{web} (a network or
  interwoven structure). In Donald Knuth's model, a program is a network
  of named sections and relationships between them, rather than merely a
  linear file.
\item
  \textbf{TANGLE} --- from the verb \emph{to tangle} (to intertwine or
  rearrange into a machine-oriented order). The utility assembles the
  compiler-oriented source code of the program.
\item
  \textbf{WEAVE} --- from the verb \emph{to weave} (to interlace or
  compose into a readable document). The utility produces a
  human-oriented document describing the program.
\end{itemize}

Technically, WEB is bilingual, and its underlying idea is designed to be
portable beyond a specific technology stack. In the original version,
TeX serves as the document-formatting language and Pascal as the
programming language, but Donald Knuth explicitly emphasises that the
principle is not tied to this pairing: TeX could be replaced with Scribe
or Troff, while Pascal could be replaced with C, LISP, FORTRAN, ALGOL,
assembly language, or other languages.

The central scheme of the section is \textbf{WEAVE / TANGLE}. WEAVE
produces a document that describes the program and facilitates its
maintenance; TANGLE produces a machine-executable program. Both results
are generated from the same source, so the code and its documentation do
not diverge into two independent versions of the truth and do not
require manual synchronisation.

This scheme also establishes the criterion for evaluating the
contemporary branches identified in the opening remarks: the primary
question is not which tools they use, but whether they preserve the
relationship between explanation, code, and a verifiable artifact.

\subsection{C--F. The Prime-Number Example: Source, Tangle, and
Weave}\label{cf.-the-prime-number-example-source-tangle-and-weave}

This section follows the methodological choice stated earlier: the
primary remaster is implemented in noweb-like C++. This format was
selected not because of any limitations in CWEB, but as a lightweight
adaptation of Donald Knuth's idea to a contemporary team workflow. Even
the basic discipline of
\texttt{canonical\ source\ -\textgreater{}\ tangling\ -\textgreater{}\ generated\ artifact\ -\textgreater{}\ smoke-check}
already imposes a non-trivial entry threshold; making a separate
readable-document generation branch mandatory would bring the example
closer to classical WEB/CWEB, but would make it less suitable for
demonstrating the minimal practical core of the method.

A more canonical contrast has been preserved in
{\urlstyle{tt}\path{02-Literate-Companion-Implementations/examples/en/01-cweb/}}:
this is a direct C-family lineage from WEB to CWEB
(\citeproc{ref-knuth1992a}{Knuth and Levy, n.d.}), with a two-branch
model:
\texttt{primes.w\ -\textgreater{}\ ctangle\ -\textgreater{}\ primes.c\ -\textgreater{}\ cc/gcc\ -\textgreater{}\ primes\ -\textgreater{}\ output.txt}
and
\texttt{primes.w\ -\textgreater{}\ cweave\ -\textgreater{}\ primes.tex\ -\textgreater{}\ pdftex\ -\textgreater{}\ primes.pdf}.
This demonstrates that the complete tangle/weave model remains
operational. For the main review, however, \texttt{02-noweb-like} was
selected in order to demonstrate the portability of the principle
without requiring entry into the CWEB/TeX toolchain.

Historically, Donald Knuth's example follows the discussion of
structured programming: the question is no longer only how to structure
control flow, but also how to structure the explanation of a program
(\citeproc{ref-dijkstra1972notes}{Dijkstra 1972: 26--39}).

The C--F subsections below therefore do not reproduce the entire
codebase. Instead, they identify the four projections of Donald Knuth's
example and immediately compare them with the corresponding stages of
our remaster pipeline. The complete noweb-like source is presented once,
in the ``Remaster'' subsection.

\subsubsection{C. Donald Knuth's Readable
Example}\label{c.-donald-knuths-readable-example}

Section C presents the readable projection of the example, which is also
discussed later in Section F: the reader is given not a linear Pascal
file, but a composition of the program as an explanation.

\begin{Shaded}
\begin{Highlighting}[]
\NormalTok{The program prints the first 1000 prime numbers:}
\NormalTok{  1. define the table parameters;}
\NormalTok{  2. introduce the table of discovered prime numbers;}
\NormalTok{  3. generate the prime numbers;}
\NormalTok{  4. print the table page by page;}
\NormalTok{  5. provide the reader with a table of contents, sections, and an index of names.}
\end{Highlighting}
\end{Shaded}

\subsubsection{D. The Original WEB File}\label{d.-the-original-web-file}

Section D presents the original \texttt{PRIMES.WEB}, that is, the
material written directly by the author: the documentation layer +
Pascal + WEB commands.

In our remaster equivalent, \texttt{primes.nw} serves as the source of
truth.

\textbf{noweb chunk markers and syntax:}

\begin{itemize}
\tightlist
\item
  \texttt{\textless{}\textless{}program\textgreater{}\textgreater{}=}
  defines a named chunk called \texttt{program}.
\item
  \texttt{\textless{}\textless{}name\textgreater{}\textgreater{}} refers
  to another named chunk.
\item
  \texttt{@} terminates the current chunk.
\item
  \texttt{notangle\ -Rprogram} selects \texttt{program} as the root
  chunk.
\item
  \texttt{notangle} expands the chunk references and generates ordinary
  C++.
\item
  The C++ compiler never sees the noweb markers; it sees only the
  generated \texttt{primes.cpp}.
\end{itemize}

\subsubsection{E. The Machine-Oriented Projection:
TANGLE}\label{e.-the-machine-oriented-projection-tangle}

Section E establishes the machine-oriented projection: \texttt{TANGLE}
transforms \texttt{PRIMES.WEB} into \texttt{PRIMES.PAS}.

In our remaster, the corresponding artifact is \texttt{primes.cpp}.

\texttt{primes.nw\ -\textgreater{}\ notangle\ -Rprogram\ -\textgreater{}\ primes.cpp\ -\textgreater{}\ g++\ -\textgreater{}\ primes\ -\textgreater{}\ output.txt}

\texttt{output.txt} serves as a verification artifact (smoke-check),
rather than as a readable document.

\subsubsection{F. The Human-Oriented Projection:
WEAVE}\label{f.-the-human-oriented-projection-weave}

Section F defines a separate branch for generating a readable document:
in WEB/CWEB, this is
\texttt{WEAVE/cweave\ -\textgreater{}\ TeX\ -\textgreater{}\ PDF}.

\subsubsection{Remaster: Why the noweb-like Version Is Deliberately
Lightweight}\label{remaster-why-the-noweb-like-version-is-deliberately-lightweight}

\texttt{02-noweb-like} presents a tangle-first C++ remaster. It
preserves the minimal practical core of the method: a canonical literate
source, named chunks, a generated machine artifact, and a smoke-check
verification artifact. A separate branch for generating a readable
document is not mandatory in this version.

The material below is presented as a unified literate form of the
remaster: the explanation remains adjacent to the corresponding chunks
rather than being extracted into separate, repetitive sections.

The top level defines the program's intent and order of composition:

\begin{Shaded}
\begin{Highlighting}[]
\NormalTok{\# Printing primes: a C++ reboot}

\NormalTok{This program prints the first 1000 prime numbers.}
\NormalTok{The top{-}level structure is deliberately simple:}
\NormalTok{generate the prime table, then print it.}

\NormalTok{\textless{}\textless{}program\textgreater{}\textgreater{}=}
\NormalTok{\#include \textless{}cstddef\textgreater{}}
\NormalTok{\#include \textless{}iomanip\textgreater{}}
\NormalTok{\#include \textless{}iostream\textgreater{}}
\NormalTok{\#include \textless{}vector\textgreater{}}

\NormalTok{\textless{}\textless{}constants\textgreater{}\textgreater{}}
\NormalTok{\textless{}\textless{}types\textgreater{}\textgreater{}}
\NormalTok{\textless{}\textless{}declarations\textgreater{}\textgreater{}}

\NormalTok{int main() \{}
\NormalTok{    const auto primes = generate\_primes(PRIME\_COUNT);}
\NormalTok{    print\_table(primes);}
\NormalTok{    return 0;}
\NormalTok{\}}

\NormalTok{\textless{}\textless{}prime{-}generation\textgreater{}\textgreater{}}
\NormalTok{\textless{}\textless{}table{-}output\textgreater{}\textgreater{}}
\NormalTok{@}
\end{Highlighting}
\end{Shaded}

The output parameters define the table layout and preserve the
connection with Donald Knuth's original example:

\begin{Shaded}
\begin{Highlighting}[]
\NormalTok{The output format follows Knuth\textquotesingle{}s original example:}
\NormalTok{50 rows per page, 4 columns, 10 character positions per column.}

\NormalTok{\textless{}\textless{}constants\textgreater{}\textgreater{}=}
\NormalTok{constexpr std::size\_t PRIME\_COUNT = 1000;}
\NormalTok{constexpr std::size\_t ROWS\_PER\_PAGE = 50;}
\NormalTok{constexpr std::size\_t COLUMNS\_PER\_PAGE = 4;}
\NormalTok{constexpr int COLUMN\_WIDTH = 10;}
\NormalTok{@}
\end{Highlighting}
\end{Shaded}

The container for the table of prime numbers is assigned a distinct
semantic role:

\begin{Shaded}
\begin{Highlighting}[]
\NormalTok{The table of primes is represented as a growing sequence.}
\NormalTok{The alias makes the role of the container explicit.}

\NormalTok{\textless{}\textless{}types\textgreater{}\textgreater{}=}
\NormalTok{using PrimeTable = std::vector\textless{}int\textgreater{};}
\NormalTok{@}
\end{Highlighting}
\end{Shaded}

The function declarations establish a separation of responsibilities
between generation, candidate validation, and output:

\begin{Shaded}
\begin{Highlighting}[]
\NormalTok{Generation, candidate checking, and output are introduced}
\NormalTok{as separate operations.}

\NormalTok{\textless{}\textless{}declarations\textgreater{}\textgreater{}=}
\NormalTok{PrimeTable generate\_primes(std::size\_t count);}
\NormalTok{bool is\_prime\_candidate(int candidate, const PrimeTable\& primes);}
\NormalTok{void print\_table(const PrimeTable\& primes);}
\NormalTok{@}
\end{Highlighting}
\end{Shaded}

Prime-number generation preserves the algorithm of the remaster version
without altering its logic:

\begin{Shaded}
\begin{Highlighting}[]
\NormalTok{\textless{}\textless{}prime{-}generation\textgreater{}\textgreater{}=}
\NormalTok{PrimeTable generate\_primes(std::size\_t count) \{}
\NormalTok{    PrimeTable primes;}
\NormalTok{    if (count == 0) return primes;}
\NormalTok{    primes.push\_back(2);}

\NormalTok{    for (int candidate = 3; primes.size() \textless{} count; candidate += 2) \{}
\NormalTok{        if (is\_prime\_candidate(candidate, primes)) \{}
\NormalTok{            primes.push\_back(candidate);}
\NormalTok{        \}}
\NormalTok{    \}}
\NormalTok{    return primes;}
\NormalTok{\}}

\NormalTok{bool is\_prime\_candidate(int candidate, const PrimeTable\& primes) \{}
\NormalTok{    for (int p : primes) \{}
\NormalTok{        if (p \textgreater{} candidate / p) return true;}
\NormalTok{        if (candidate \% p == 0) return false;}
\NormalTok{    \}}
\NormalTok{    return true;}
\NormalTok{\}}
\NormalTok{@}
\end{Highlighting}
\end{Shaded}

The output phase preserves the page-by-page layout and the calculation
of indices by row and column:

\begin{Shaded}
\begin{Highlighting}[]
\NormalTok{\textless{}\textless{}table{-}output\textgreater{}\textgreater{}=}
\NormalTok{void print\_table(const PrimeTable\& primes) \{}
\NormalTok{    std::size\_t page\_number = 1;}
\NormalTok{    std::size\_t page\_offset = 0;}

\NormalTok{    while (page\_offset \textless{} primes.size()) \{}
\NormalTok{        std::cout \textless{}\textless{} "The First " \textless{}\textless{} primes.size()}
\NormalTok{                  \textless{}\textless{} " Prime Numbers {-}{-}{-} Page " \textless{}\textless{} page\_number \textless{}\textless{} "\textbackslash{}n\textbackslash{}n";}

\NormalTok{        for (std::size\_t row = 0; row \textless{} ROWS\_PER\_PAGE; ++row) \{}
\NormalTok{            for (std::size\_t column = 0; column \textless{} COLUMNS\_PER\_PAGE; ++column) \{}
\NormalTok{                const std::size\_t index = page\_offset + row + column * ROWS\_PER\_PAGE;}
\NormalTok{                if (index \textless{} primes.size()) \{}
\NormalTok{                    std::cout \textless{}\textless{} std::setw(COLUMN\_WIDTH) \textless{}\textless{} primes[index];}
\NormalTok{                \}}
\NormalTok{            \}}
\NormalTok{            std::cout \textless{}\textless{} \textquotesingle{}\textbackslash{}n\textquotesingle{};}
\NormalTok{        \}}

\NormalTok{        std::cout \textless{}\textless{} \textquotesingle{}\textbackslash{}n\textquotesingle{};}
\NormalTok{        ++page\_number;}
\NormalTok{        page\_offset += ROWS\_PER\_PAGE * COLUMNS\_PER\_PAGE;}
\NormalTok{    \}}
\NormalTok{\}}
\NormalTok{@}
\end{Highlighting}
\end{Shaded}

\texttt{primes.nw} remains the canonical literate source;
\texttt{primes.cpp} serves as the generated machine artifact; and
\texttt{output.txt} serves as the smoke-check output.

\subsection{G. Additional Bells and
Whistles}\label{g.-additional-bells-and-whistles}

In Section G, Donald Knuth makes a transition: the prime-number example
is useful as a demonstration, but it does not establish WEB's
suitability for real-world development. A system must be capable of
handling not only an instructional fragment, but also large programs,
different compilers, non-standard conventions, typography, portability,
and the accumulation of technical exceptions. Otherwise, it is not a
method, but merely an elegant toy constructed for an article.

Donald Knuth lists several capabilities of WEB83 that were not required
in \texttt{PRIMES.WEB}, but become necessary in larger programs: manual
control over WEAVE formatting, the ability to pass fragments through
TANGLE almost verbatim, control over spacing and line beginnings, octal
and hexadecimal constants, ASCII-based character encoding, a string pool
with a checksum, numeric macros, and simple compile-time arithmetic.

The significance of this section does not lie in the features
themselves. Many of them now appear to be compensations for the
limitations of Pascal and older compilers. In C++, some of these
problems simply disappear: the language provides proper strings,
\texttt{constexpr}, local variables, standard containers, formatters,
and build systems.

The point is that a literate tool must be capable of working with the
messy realities of a language and platform, rather than only with an
idealised example. The prime-number example explains the principle, not
the entire production toolchain.

For the noweb-like C++ remaster, this can be expressed as follows:

\begin{Shaded}
\begin{Highlighting}[]
\NormalTok{Toy example:}
\NormalTok{  \textless{}\textless{}program\textgreater{}\textgreater{}}
\NormalTok{  \textless{}\textless{}constants\textgreater{}\textgreater{}}
\NormalTok{  \textless{}\textless{}types\textgreater{}\textgreater{}}
\NormalTok{  \textless{}\textless{}prime{-}generation\textgreater{}\textgreater{}}
\NormalTok{  \textless{}\textless{}table{-}output\textgreater{}\textgreater{}}

\NormalTok{Real{-}world program:}
\NormalTok{  \textless{}\textless{}headers\textgreater{}\textgreater{}}
\NormalTok{  \textless{}\textless{}platform{-}specific{-}code\textgreater{}\textgreater{}}
\NormalTok{  \textless{}\textless{}configuration\textgreater{}\textgreater{}}
\NormalTok{  \textless{}\textless{}generated{-}bindings\textgreater{}\textgreater{}}
\NormalTok{  \textless{}\textless{}tests\textgreater{}\textgreater{}}
\NormalTok{  \textless{}\textless{}diagnostics\textgreater{}\textgreater{}}
\NormalTok{  \textless{}\textless{}build{-}notes\textgreater{}\textgreater{}}
\end{Highlighting}
\end{Shaded}

Thus, the ``additional bells and whistles'' are not embellishments in
the ordinary sense. They form a layer that allows the literate source to
remain the principal source even when the project ceases to be clean and
small. The circumstances in which such discipline is genuinely justified
will be examined in the conclusion of this review; Donald Knuth himself
also proceeds to qualify the universality of WEB and does not present it
as a tool for every possible case.

\subsection{H. Occam's Razor}\label{h.-occams-razor}

Section H serves as a methodological qualification of the preceding
section. After a lengthy list of capabilities, it is easy to draw the
false conclusion that a literate system should continue expanding
indefinitely and incorporate everything: formatting, portability, macro
expansion, publication, build processes, diagnostics, and environment
management. Donald Knuth instead invokes Occam's razor: complexity
should be introduced only where it is genuinely required by the task,
rather than merely because the tool is technically capable of
accommodating yet another feature.

This is particularly important in considering the history of WEB, whose
original strength lay in connecting the explanation and the code of a
program within a single working process. Yet the same system
historically accumulated ceremonial overhead and a substantial entry
threshold: numerous commands, a typographic layer, and a dependence on
the conventions of a particular school of development. Later branches
therefore embodied a different response: not the indefinite expansion of
the core, but its reduction to a minimally viable discipline.

In this sense, noweb, as used in our remaster, is significant as an
engineering reduction rather than as an ``anti-WEB''. It preserves the
essential relationship between named fragments and tangling, while
removing much of the ceremonial overhead. This neither supersedes
classical WEB nor makes noweb a universal solution; it demonstrates that
a literate approach requires proportion. If a system attempts to
encompass every task simultaneously, it begins to succumb to its own
complexity before it can systematically improve understanding.

\subsection{I. Portability}\label{i.-portability}

In Section I, Donald Knuth moves from clarity of exposition to
portability. TeX and its associated software had to operate across
numerous machines and operating systems. WEB proved useful precisely
because it made it possible to preserve a shared program source while
still accounting for platform-specific differences.

Donald Knuth first describes a bootstrapping scheme: given
\texttt{TANGLE.PAS}, \texttt{TANGLE.WEB}, and \texttt{WEAVE.WEB}, one
can obtain a working version of TANGLE, use it to generate
\texttt{WEAVE.PAS}, and then build the entire WEB system and programs
such as TeX.

Donald Knuth then qualifies the practical limitations of the model. Real
machines differ in character sets, file conventions, terminals, data
packing, floating-point arithmetic, and compiler quality. It is
therefore methodologically unsound to expect a single
\texttt{TANGLE.PAS} or \texttt{TANGLE.WEB} to operate everywhere without
modification. System-dependent changes are inevitable.

WEB's solution is \textbf{change files}. TANGLE and WEAVE read not only
the primary WEB file, but also a separate \texttt{.CH} file containing
substitutions for a particular system. As a result, the master source
remains stable, while platform-specific modifications are maintained
separately. Donald Knuth emphasises that \texttt{TANGLE.WEB} itself
remains effectively unchanged: it is \texttt{TANGLE.CH} that varies,
while the core logic of TANGLE remains shared across systems.

Donald Knuth's approach to literate programming is not detached from
build and distribution processes.

For our noweb-like C++ remaster, this can be represented as follows:

\begin{Shaded}
\begin{Highlighting}[]
\NormalTok{canonical source:}
\NormalTok{  primes.nw}

\NormalTok{generated artifact:}
\NormalTok{  primes.cpp}

\NormalTok{platform overlay:}
\NormalTok{  primes.linux.patch}
\NormalTok{  primes.windows.patch}
\NormalTok{  config/platform.hpp}
\NormalTok{  CMake presets}
\end{Highlighting}
\end{Shaded}

In other words, a direct contemporary analogue of .CH does not need to
be a separate ``change file'' in the classical WEB format. In modern
practice, the range of productive analogies is broader:

\begin{Shaded}
\begin{Highlighting}[]
\NormalTok{change file → patch layer}
\NormalTok{change file → platform adapter}
\NormalTok{change file → environment{-}specific config}
\NormalTok{change file → build preset}
\NormalTok{change file → overlay}
\end{Highlighting}
\end{Shaded}

The principle, however, remains the same: the primary explanatory source
should not be disrupted by every platform-specific requirement. If
Windows, Linux, or a particular compiler requires divergent behaviour,
that divergence should be isolated within an adaptation layer.
Otherwise, the canonical source rapidly becomes a collection of
platform-specific exceptions in which the central idea is obscured by
adaptation noise.

Contemporary branches implement this principle through different means,
ranging from master-source discipline to the management of reproducible
environments.

The distinction between classical WEB and the executable-document branch
becomes substantive at this point. For Donald Knuth, portability
primarily concerns transferring a program between machines. In
Quarto/Jupyter/R Markdown, portability often means something different:
the ability to rebuild a computational report with the same
dependencies, the same data, and the same execution order. This is not a
literal \texttt{.CH}, but it reflects the same underlying concern: the
source must survive a change of environment.

Contemporary development introduces an additional layer: portability is
no longer limited to differences between operating systems and
compilers. Increasingly, the issue is a reproducible working
environment: WSL\footnote{Windows Subsystem for Linux, or WSL, is a
  Windows feature that runs Linux environments and Linux tools without
  requiring the separate setup of a virtual machine or a dual-boot
  configuration.}, containers\footnote{A container is an isolated
  runtime environment in which an application runs with explicitly
  specified dependencies and configuration.}, cloud IDEs\footnote{A
  cloud IDE is a development environment delivered through a browser or
  remote infrastructure rather than only through the author's local
  machine.}, CI/CD runners\footnote{A CI/CD runner is the executor of an
  automated pipeline for checking, building, or publishing a project.},
dev containers\footnote{A dev container describes a reproducible
  development environment, including tools, extensions, dependencies,
  and project commands.}, and managed environments. Under these
conditions, the question changes from ``will the program build on
another machine?'' to ``can the entire process be restored: source,
tooling, dependencies, build commands, verification artifacts, and
execution order?''

It is illustrative that the companion examples for this review were
reproduced within a WSL environment under Windows 11. This is a
contemporary form of portability: Windows serves as the working
platform, WSL as the Linux execution environment, and the examples
themselves are verified using specific CLI tools\footnote{A command-line
  interface, or CLI, controls a tool through terminal commands, which
  makes checks and automation easier to reproduce.}: \texttt{ctangle},
\texttt{notangle}, \texttt{emacs}, \texttt{quarto}, \texttt{jupyter},
\texttt{Rscript}, \texttt{pandoc}, and \texttt{g++}. Portability thus
becomes not an abstract property of the source text, but a verifiable
state of the entire toolchain.

The changing role of the environment is particularly clear. In classical
WEB, system-dependent differences were isolated in \texttt{.CH} files;
in contemporary practice, an analogous function may be performed by a
WSL profile, a container, a lockfile\footnote{A lockfile records exact
  dependency versions and related metadata so that repeated installation
  is not affected by incidental upstream updates.}, a CI configuration,
a build preset\footnote{A build preset is a named build profile with
  predefined parameters such as compiler, mode, and output directory.},
or reproducible-run documentation. The principle remains unchanged: the
primary explanatory source should not disintegrate because of the
peculiarities of a particular machine. What must now be recorded
alongside the source is not only the platform-specific modification, but
also the means of reconstructing the execution environment itself.

\subsection{J. Programs as Webs}\label{j.-programs-as-webs}

Section J returns to the article's central architectural idea: a program
is structured neither as a simple sequence of steps nor as a tidy
top-down tree. Donald Knuth proposes that it should instead be
understood as a network --- a web --- in which named fragments are
connected by meaning, dependencies, and explanatory order. One node may
introduce an idea, another refine its constraints, and a third provide
its implementation; the reader moves through this network deliberately,
rather than being compelled to follow the order imposed by the compiler.

\begin{quote}
Retrospective: This section appears methodologically farsighted. It is
not a direct prediction of contemporary tools, but Donald Knuth's model
maps closely onto later problems of navigation within complex codebases:
dependency graphs, module maps, symbol navigation in IDEs,
call/reference navigation, knowledge graphs, and project context maps;
LLM context maps can be mentioned here only as a preliminary horizon.
For contemporary practice, the implication is straightforward:
understanding a program is increasingly constructed through a network of
relationships between fragments, rather than through the linear
inspection of files.
\end{quote}

The WEB model cannot be reduced to the attractive metaphor that
``everything is connected to everything else''. Its value lies in the
fact that the network of relationships is recorded in a verifiable
source from which both the machine-oriented result and the
human-oriented explanation are generated. This is why Section J occupies
a position between portability and style in the logic of our review:
first, the integrity of the source is preserved across platforms; then,
the source is understood as a network; only after that do we consider
how this network should be named and presented.

\subsection{K. Stylistic Issues}\label{k.-stylistic-issues}

In Section K, Donald Knuth moves from the architecture of WEB to the
style of working within it. Once a program becomes an explanatory text,
style ceases to be cosmetic: it becomes part of the engineering of
understanding.

The central question of the section is how sections should be named and
organised. Donald Knuth distinguishes between macros and named
fragments: small technical abbreviations may remain macros, but larger
parts of a program should receive human-readable names. Such a name need
not be formally exhaustive; it should express the meaning of the
fragment with sufficient precision, without burdening the reader with
excessive detail.

A poor section name in the noweb-like C++ remaster:

\begin{Shaded}
\begin{Highlighting}[]
\NormalTok{\textless{}\textless{}code{-}for{-}loop\textgreater{}\textgreater{}}
\end{Highlighting}
\end{Shaded}

A better name:

\begin{Shaded}
\begin{Highlighting}[]
\NormalTok{\textless{}\textless{}generate prime candidates\textgreater{}\textgreater{}}
\end{Highlighting}
\end{Shaded}

Better still, if the fragment genuinely expresses an action:

\begin{Shaded}
\begin{Highlighting}[]
\NormalTok{\textless{}\textless{}print one page of the prime table\textgreater{}\textgreater{}}
\end{Highlighting}
\end{Shaded}

Donald Knuth emphasises the balance between formal and informal
exposition. A section name should not be turned into a complete
specification of every variable and precondition: if it is, it ceases to
assist the reader and begins to imitate a legal contract. Yet it must
not remain excessively general either. A well-named section should be
precise enough for the reader to understand its role before reading the
code.

For the C++ version, this means that fragments should be named according
to their purpose rather than their syntax:

\begin{Shaded}
\begin{Highlighting}[]
\NormalTok{\textless{}\textless{}constants\textgreater{}\textgreater{}}
\NormalTok{\textless{}\textless{}types\textgreater{}\textgreater{}}
\NormalTok{\textless{}\textless{}declarations\textgreater{}\textgreater{}}
\end{Highlighting}
\end{Shaded}

--- are acceptable for auxiliary blocks, but semantic names are
preferable for logic:

\begin{Shaded}
\begin{Highlighting}[]
\NormalTok{\textless{}\textless{}generate the requested number of primes\textgreater{}\textgreater{}}
\NormalTok{\textless{}\textless{}test whether a candidate is prime\textgreater{}\textgreater{}}
\NormalTok{\textless{}\textless{}print the prime table page by page\textgreater{}\textgreater{}}
\end{Highlighting}
\end{Shaded}

Donald Knuth also discusses control structures separately. If the
control flow is non-standard, the section name should make this
explicit. The danger lies not in a \texttt{goto}, loop, or early exit as
such, but in an unexplained control transition. In a contemporary C++
version, the same concern applies to \texttt{break}, \texttt{continue},
\texttt{return}, exceptions, guard clauses, and other means of altering
the ordinary flow of a function.

Across all branches, style remains a working instrument: the easier it
is to navigate the names and roles of fragments, the more robust the
understanding of the program becomes.

Later engineering culture develops a similar discipline through naming
conventions: the names of classes, methods, APIs, modules, UI blocks,
and BEM-like schemes\footnote{BEM, or Block--Element--Modifier, is a
  UI-component naming convention in which a name indicates the block,
  its element, and a state or modification variant.}, as codified in
style guides and readable-code practices. This is not literate
programming in itself, but it reflects the same effort to construct a
readable structure: a name should communicate the role of a fragment
before the reader examines its implementation.

This section therefore draws a boundary between literate programming and
contemporary documentation generators. Javadoc, Doxygen, or comment
extraction can produce an API reference, but this is not necessarily
literate programming. Such tools more often describe an already existing
code structure. In Donald Knuth's approach, style operates
prospectively.

\subsection{L. Economic Issues}\label{l.-economic-issues}

In Section L, Donald Knuth asks a pragmatic question: what does WEB
cost? He begins by considering direct computational expenses. TANGLE
takes approximately as much time as compiling the resulting Pascal file;
WEAVE also runs reasonably quickly, although TeX typesetting is already
noticeably more expensive. It is therefore not always sensible to
rebuild and print the documentation in full several times a day.

Documentation is rarely printed today, but the central economic argument
does not concern processor time or the cost of paper. Donald Knuth
argues that the total time he spends writing and debugging a WEB program
is no greater than the time required to write and debug an ordinary
ALGOL or Pascal program, despite the substantially more extensive
documentation that results. The additional time devoted to explanation
is recovered through reduced debugging effort, because the mode of
exposition forces the author to clarify ideas earlier. When a program is
written ``for oneself'', it is easy to rely on shortcuts that later
become sources of error. When a program is written as an explanation, it
becomes more difficult for the author to deceive himself.

For the noweb-like C++ remaster, this means that the cost of the method
must be stated plainly:

\begin{Shaded}
\begin{Highlighting}[]
\NormalTok{a higher entry threshold;}
\NormalTok{the need for a tangle/render workflow;}
\NormalTok{the need to maintain canonical{-}source discipline;}
\NormalTok{the requirement not to edit the generated .cpp manually;}
\NormalTok{the need to maintain the explanation alongside the code.}
\end{Highlighting}
\end{Shaded}

The benefits are equally concrete:

\begin{Shaded}
\begin{Highlighting}[]
\NormalTok{intent is recorded more clearly;}
\NormalTok{review becomes easier;}
\NormalTok{returning to the code later becomes easier;}
\NormalTok{introducing a new reader becomes easier;}
\NormalTok{code, tests, and documentation are easier to connect.}
\end{Highlighting}
\end{Shaded}

The economics of the different branches vary in their entry thresholds
and risk profiles, but the underlying exchange follows the same model:
greater discipline at the outset in return for lower maintenance costs
later.

This section also provides a bridge to contemporary reproducibility
practices. In research and data science, the literate approach became
economically justified because reports, code, graphs, results, and
validation began to be assembled from a single source. In widespread
practice, this is visible in executable documents and reproducible
computation: the idea has become part of the everyday workflow of
science and data science.

The section concludes with an irony: Donald Knuth worries that literate
programs may appear so complete to their authors that they will feel
compelled to publish them everywhere. Once code becomes text, the
temptation to aestheticise it inevitably emerges.

For the purposes of this review, the boundary must remain clear: the
value lies not in making a program resemble a ``work of art'', but in
making it better explained and more readily verifiable.

\subsection{M. Related Work}\label{m.-related-work}

In Section M, Donald Knuth removes the aura of solitary invention from
WEB. He states explicitly that the system contains nothing ``really
new'': he assembled ideas that had existed for a long time and combined
them into a working form. This is an important act of intellectual
candour: WEB is presented as a node in the history of programming,
typography, and the publication of algorithms.

The first important source is George Forsythe, for whom a useful
algorithm constitutes a contribution to knowledge, while its publication
is a form of scholarly work. This directly supports Donald Knuth's
argument: a program can be not merely a functional mechanism, but also a
publishable intellectual object. The idea of a program as literature
follows naturally from this position, but without romanticisation: the
issue is the scholarly form in which an algorithm is presented.

Donald Knuth then refers to Pierre-Arnoul de Marneffe and holon
programming, followed by Dijkstra, Hoare, Dahl, Wirth, and Naur. The
central intellectual lineage becomes visible here: abstraction,
structured programming, a balance between formal and informal
description, and the publication of well-written programs. WEB does not
supersede this tradition; it attempts to give it an instrumental form.

Particularly important is the episode involving Tony Hoare, who
encouraged Donald Knuth to publish TeX. For Knuth, this posed a
challenge: it is one thing to publish polished toy problems, and quite
another to expose a substantial real-world program with all its
compromises, portability concerns, and untidy details. This problem
generated the practical need for WEB: to make a large program accessible
to public reading.

The section shows that literate programming emerged at the intersection
of several lineages:

\begin{Shaded}
\begin{Highlighting}[]
\NormalTok{structured programming;}
\NormalTok{the publication of algorithms;}
\NormalTok{the typography of programs;}
\NormalTok{documentation as part of design;}
\NormalTok{a balance between formal and informal exposition;}
\NormalTok{the attempt to make a large program readable.}
\end{Highlighting}
\end{Shaded}

The contemporary landscape of related practices has a similar structure.
CWEB, noweb, Org Babel, R Markdown, Quarto, and Jupyter do not form a
single direct lineage. They are better understood as a family of
practices in which different domains address related engineering
problems in different ways. This landscape can be read in terms of a
two-layer influence: direct or historically proximate continuations of
WEB --- above all, CWEB and noweb --- and the more widespread layer of
executable documents and reproducible computing, which addresses related
problems through different means.

The conclusion should therefore avoid searching for a single ``true
successor''. CWEB is closest genealogically; noweb, in the simplicity of
its mechanism; Org Babel, in its plain-text practice; and
Quarto/Jupyter/R Markdown, in their institutional scale. Each implements
part of the same engineering task in its own way. What matters is that
direct continuation of WEB, instrumental affinity, and productive
conceptual association should not be conflated.

Ultimately, WEB is a synthesis rather than a solitary invention. This,
in turn, legitimises the comparative approach: contemporary branches
should likewise be read not as copies of WEB, but as different ways of
continuing its central inquiry --- how to make a program readable and
verifiable.

\subsection{N. Retrospect and
Prospects}\label{n.-retrospect-and-prospects}

In the final section, Donald Knuth himself tempers the rhetoric of his
manifesto. He acknowledges that the creators of new languages are
generally inclined to overestimate the significance of their own
experience, and that his work with WEB is too deeply shaped by his
personal preferences.

The article ends not with a declaration of victory, but with a cautious
test: will the method work beyond its author and his environment?

At the same time, Donald Knuth formulates a strong conclusion: in his
experience, WEB makes it possible to create programs that are more
portable, more comprehensible, and easier to maintain, while the method
works not only for small examples but also for large systems. For him,
WEB is the result of a reversal in the relationship between typography
and computing: the computer was first applied to typesetting, and
typography then returned to the centre of computer science as a means of
making programs readable.

Donald Knuth then explicitly limits the intended audience. WEB was not
designed for everyone. Its user must be prepared to work with several
languages simultaneously: WEB, TeX, Pascal, and the algorithmic errors
of the program itself. It is not a low-friction tool, but an environment
for people who enjoy writing and explaining what they do. This helps to
explain why classical WEB did not become a mass standard for ordinary
software development.

Yet Donald Knuth's forecast proved more accurate than it might initially
appear. He anticipates a future in which the philosophy of WEB is
embedded in more efficient programming environments: tangling and
compiling might be combined, debugging might operate in terms of the WEB
source, and program sections might be displayed on demand rather than
requiring the entire document to be printed. This already resembles
IDEs, symbol navigation, selective rendering, notebook environments, and
contemporary context-management tools.

These lines of development can then be brought together in a two-track
conclusion: WEB-like source generation and executable-document
narrative.

The final section thus transforms the particular experience of WEB into
a broader conclusion: classical WEB remained a niche tool, but its
principles spread far more widely.

\subsection{Conclusion. 1. What Donald Knuth Actually
Proposed}\label{conclusion.-1.-what-donald-knuth-actually-proposed}

Donald Knuth's central proposal cannot be reduced to improved
commenting. A comment usually remains a secondary layer: it explains
code that has already been organised in an order determined by the
compiler, the programming language, or the author's habits.
\emph{Literate programming} proposes a stronger model: a program should
be constructed from the outset as an explanatory document intended for
human reading.

In this model, the central object is not the generated \texttt{.pas},
\texttt{.c}, \texttt{.cpp}, or \texttt{.html} file, but a unified source
from which different representations are derived:

\begin{Shaded}
\begin{Highlighting}[]
\NormalTok{literate source}
\NormalTok{  → machine{-}oriented artifact}
\NormalTok{  → human{-}oriented document}
\end{Highlighting}
\end{Shaded}

In Donald Knuth's implementation, this scheme is realised through WEB,
TeX, Pascal, TANGLE, and WEAVE, but the underlying principle is broader
than the original technological pairing. What matters is not the
specific historical toolchain, but the discipline: explanation, code,
and the mechanism for producing artifacts must remain connected. This is
precisely why a program in Donald Knuth's model ceases to be a linear
file. It becomes a network of named fragments in which the order of
exposition can follow the logic of understanding rather than the
sequence required by the machine.

Code is part of the explanatory structure; this remains the principal
criterion for evaluating contemporary continuations and
reinterpretations of literate programming.

\subsection{Conclusion. 2. What Persisted and What Changed
Form}\label{conclusion.-2.-what-persisted-and-what-changed-form}

Classical WEB did not become a mass standard for software development,
but several of its principles proved durable. They continued to exist
not within a single lineage, but across different engineering and
research practices.

The first lineage is \textbf{source generation}. CWEB, noweb, and Org
Babel preserve the idea closest to WEB: a structured source exists from
which machine-oriented code can be produced. Named fragments, tangling,
and the prohibition against manually editing derived artifacts are
particularly important within this lineage.

The second lineage is \textbf{executable documents}. Quarto, Jupyter,
and R Markdown shift the emphasis from generating program source code to
producing a reproducible document in which text, code, execution, and
results belong to a single process. This is no longer WEB in the literal
sense, but the task remains related: a computation must not only be
performed, but also explained, rebuilt, and verified.

The third lineage is \textbf{navigation within complex systems}. Donald
Knuth's section on programs as webs proved particularly farsighted.
Contemporary IDEs, dependency graphs, symbol navigation, module maps,
knowledge graphs, and context maps address a similar problem: a program
is understood not through the linear reading of files, but through a
network of relationships between fragments.

The fourth lineage is \textbf{the reproducible environment}. In
classical WEB, portability was supported through a master source and
change files. In contemporary practice, WSL, containers, CI, dev
containers, lockfiles, build presets, and documented execution commands
partially fulfil the same role. The companion suite for this review was
verified within a WSL environment under Windows 11, which illustrates
the new form of an old question: what must be reproduced is not only the
code, but also the environment in which the code becomes a verifiable
artifact.

It is therefore more accurate to speak not of a single successor to WEB,
but of a distributed landscape of continuations, reinterpretations, and
conceptual associations. \texttt{01-cweb}, \texttt{02-noweb-like}, and
\texttt{03-org-babel} form the WEB-like / source-generation lineage.
\texttt{04-quarto}, \texttt{05-jupyter}, and \texttt{06-rmarkdown} form
the executable-document / computational-publishing lineage. These
branches differ technically, but each continues Donald Knuth's central
inquiry in its own way: how to connect a program, its explanation, and a
verifiable result.

\subsection{Conclusion. 3. The Cost of the Method and Its Limits of
Applicability}\label{conclusion.-3.-the-cost-of-the-method-and-its-limits-of-applicability}

Literate programming is not a universal replacement for ordinary
software development. Its strength becomes apparent where understanding
the program is itself part of the production result: in complex
algorithms, research computing, systems with long life cycles,
domain-specific rules, generators, DSLs, architectural pipelines,
educational materials, and projects in which the cost of an error
exceeds the cost of explaining the system in advance.

Its limitations are also substantial. Classical WEB requires proficiency
in several languages and modes of work: the document language, the
programming language, the syntax of the literate system, and the domain
problem itself. Even the lightweight noweb-like approach retains the
cost of discipline: the canonical source must be maintained,
tangling/rendering must be performed, generated artifacts must not be
edited manually, and the result must be checked regularly.

This cost is not a defect of the method; it defines its scope of
applicability. For short scripts, one-off utilities, and trivial CRUD
code, a full literate process will often be excessive. For systems that
must be explained, maintained, verified, and transferred to other
people, it becomes an engineering means of reducing the future cost of
misunderstanding.

A cautious connection with the language of engineering patterns is
appropriate here. Literate programming is not a GoF pattern in the
strict catalogue sense (\citeproc{ref-gamma1994designPatterns}{Gamma et
al. 1994}). It can, however, be understood as a durable pattern for
organising source and thought: a recurring solution to a problem in
which code, explanation, and verification must remain connected. In this
sense, it is closer not to a specific object-oriented template, but to a
methodological framework for design.

This concludes the retrospective part of the review proper. The scope
will now be deliberately expanded: the discussion will concern not only
the forms that literate programming has historically taken and the
lineages that can be traced through contemporary tools, but also the
ways in which Donald Knuth's central principle may be reinterpreted
under the conditions of LLM-assisted development. This is not a claim of
direct historical succession, but an authorial development of the
problem posed by Knuth.

\subsection{Conclusion. 4. LLMs: A Continuation of the Idea or a New
Form of
Rupture}\label{conclusion.-4.-llms-a-continuation-of-the-idea-or-a-new-form-of-rupture}

Large language models (LLMs) make Donald Knuth's question newly urgent.
If the source of truth is reduced to a prompt alone, literate
programming does not emerge. What results is the rapid generation of
code from an intention, but not necessarily an explainable, verifiable,
and maintainable source. A prompt may be a useful input, but by itself
it rarely records the architecture, constraints, invariants, test
scenarios, and traceability of the accepted decision.

The difference between TANGLE and an LLM is fundamental. TANGLE
deterministically transforms a structured source into a program
artifact. An LLM probabilistically continues text on the basis of
context, instructions, and examples. LLM output therefore cannot be
treated as a build branch: it must be regarded as a candidate artifact
that undergoes review and testing before being incorporated back into a
managed source.

A strong contemporary version looks as follows:

\begin{Shaded}
\begin{Highlighting}[]
\NormalTok{human{-}authored literate plan}
\NormalTok{  → chunk contracts}
\NormalTok{  → bounded prompt}
\NormalTok{  → LLM{-}generated candidate}
\NormalTok{  → review}
\NormalTok{  → tests / smoke{-}check}
\NormalTok{  → TRACE}
\NormalTok{  → accepted update to canonical source}
\end{Highlighting}
\end{Shaded}

In this model, the LLM is not the architect by default. It operates
within a structure defined by a human: sections, chunks, constraints,
expected behaviour, and acceptance criteria must exist before
generation, or must be refined before the next run. Otherwise, the LLM
does not continue literate programming; it merely accelerates the
production of code that must subsequently be understood from scratch.

Adjacent practices --- scripting, dialogue systems, rule engines, query
parsers, domain-specific languages (DSLs)\footnote{A domain-specific
  language, or DSL, is a language deliberately restricted to the tasks
  of a particular domain.}, and behaviour trees --- help reveal the
shared engineering motive. An intermediate language of intentions often
appears between the human and the machine. In a DSL or rule engine,
however, this language is usually formalised, whereas an LLM introduces
probabilistic interpretation. For AI-assisted development, the
discipline of source, contracts, tests, and traceability therefore
becomes more important.

This is precisely where the boundary with vibe coding lies. A simple
prompt → code scheme may be productive for drafting, but it does not in
itself continue literate programming. A substantive continuation is
possible only within a workflow in which the prompt is subordinate to an
explanatory source, while the result is verified and returned to a
traceable project structure.

This is why the companion suite contains no deterministic
\texttt{07-llm} equivalent to CWEB, noweb, Org Babel, Quarto, Jupyter,
or R Markdown. Instead, \texttt{07-prompt-literate} has been introduced
as a demonstrative application of \textbf{Prompt-Literate Workflow}. Its
purpose is not to reproduce LLM output bit for bit, but to demonstrate a
controlled process: a human-authored plan, chunk contracts, bounded
prompts, review, smoke-check, and TRACE.

\texttt{07-prompt-literate} does not redefine Prompt-Literate Workflow.
It is a local example of applying the method within the article to
Donald Knuth's prime-number problem. Result-specific markers and local
checks are implemented as an additive extension of the core methodology:
it introduces local verification rules without overriding the basic
invariants.

Industrial tools are already moving in this direction: VS Code / GitHub
Copilot custom instructions, prompt files, agent customisation,
skills/hooks, and MCP-like mechanisms\footnote{Model Context Protocol,
  or MCP, is an open protocol for connecting AI applications to external
  systems, data sources, and tools.} establish a persistent project
context (\citeproc{ref-vscodeCopilotCustomInstructions}{{``Customize
GitHub Copilot in Visual Studio Code,''} n.d.};
\citeproc{ref-githubCopilotCustomization}{{``Customizing GitHub
Copilot,''} n.d.}). Regulatory requirements and provenance
standards\footnote{Provenance denotes information about an artifact's
  origin, authorship, and chain of modifications.} provide an additional
context: the EU AI Act\footnote{The EU AI Act is the European Union
  regulation that sets risk-based requirements for
  artificial-intelligence systems.} and the General-Purpose AI Code of
Practice\footnote{The General-Purpose AI Code of Practice is a voluntary
  EU framework designed to help providers of general-purpose AI models
  comply with the requirements of the EU AI Act.} intensify the demand
for transparency and accountability, while C2PA\footnote{C2PA is a
  technical specification developed by the Coalition for Content
  Provenance and Authenticity for machine-readable records of
  digital-content origin and modification history.} provides a technical
model for recording the provenance and modification history of digital
content. These remain indicators of transitional standardisation,
however, rather than conclusive evidence that AI-assisted coding has
already become literate programming.

\subsection{Conclusion. 5. Where the Idea May Develop
Further}\label{conclusion.-5.-where-the-idea-may-develop-further}

The most evident area of application for the literate approach today is
scientific and analytical documentation. In this domain, the connection
between text, code, and computation --- a connection closely related to
Donald Knuth's original formulation --- has already become an everyday
practice: R Markdown, Quarto, knitr, Org Babel, and notebook
environments connect text, code, data, computations, graphs, and
publication. In these fields, explanation is not an embellishment; it is
necessary for verifying the result.

The second area is education. The literate approach is particularly
useful where it is necessary to show not only the final code, but also
the reasoning behind it: algorithms, data structures, systems
programming, C++, and the analysis of legacy code. A program becomes not
a ``correct answer'' appended to the end of educational material, but a
route towards understanding.

The third area is architectural and tooling documentation: build
pipelines, code generation, DSLs, configuration systems, runtime
contracts, test scenarios, and reproducible environments. Conventional
documentation placed alongside the code quickly becomes outdated if it
is not connected to artifacts and verification procedures. A literate
source does not solve this problem automatically, but it compels the
author to define the source of truth explicitly.

The fourth area is game development and game design document (GDD)
pipelines\footnote{A game design document, or GDD, is a project document
  that connects a game's concept, rules, content, systems, and future
  implementation.}. Here, Donald Knuth's idea can be extended beyond
programming in the strict sense. A GDD already occupies an intermediate
position between literary description, specification, a system of rules,
and future implementation. A good GDD does not merely describe a game:
it should generate scenes, mechanics, data, quest structures, test
scenarios, tasks, and verifiable artifacts.

At this point, a broader framework emerges: \textbf{literate design}.

\begin{Shaded}
\begin{Highlighting}[]
\NormalTok{GDD source}
\NormalTok{  → narrative documentation}
\NormalTok{  → mechanics/spec chunks}
\NormalTok{  → quest structures}
\NormalTok{  → data schemas}
\NormalTok{  → implementation tasks}
\NormalTok{  → validation checks}
\NormalTok{  → generated wiki / build artifacts}
\end{Highlighting}
\end{Shaded}

This is not a literal continuation of WEB, but it extends the original
question into another domain: how can a complex system be described in a
form from which working representations and verifiable artifacts can be
produced?

For our projects, this review has a practical purpose. It did not
originate as a nostalgic commentary on an article published in 1984, but
as preparation for a pipeline in which a unified explanatory document
can connect a GDD, code, tests, a wiki, generated artifacts, and
LLM-assisted work. Our source must be sufficiently expressive to explain
the system and sufficiently disciplined to generate verifiable results.

Practical application has shown that a universal workflow must remain
compact. Domain-specific rules --- the architecture of a particular
project, runtime integration, persistence, authoring tools, and the
publication pipeline --- should be introduced as project-local
extensions and must under no circumstances become part of the core
method.

Literate programming is not required everywhere. But where understanding
is itself part of production, its central idea remains relevant: first
construct an explanatory source, then derive machine-oriented and
human-oriented representations from it, rather than attempting to
reconstruct meaning after the code has already emerged.

\protect\phantomsection\label{refs}
\subsection*{A. Primary Knuth / WEB / CWEB, structured-programming, and software-design context}
\begin{CSLReferences}{1}{0}
\hypertarget{ref-dijkstra1972notes}{}
\bibitem[\citeproctext]{ref-dijkstra1972notes}
Dijkstra, E.W. 1972: {`Notes on Structured Programming'}. In O.-J. Dahl
et al. (eds), \emph{Structured Programming}. Academic Press, 1--82.

\hypertarget{ref-gamma1994designPatterns}{}

\bibitem[\citeproctext]{ref-gamma1994designPatterns}
Gamma, E., Helm, R., Johnson, R., and Vlissides, J. 1994: \emph{Design
Patterns: Elements of Reusable Object-Oriented Software}.
Addison-Wesley.

\hypertarget{ref-knuth1984a}{}

\bibitem[\citeproctext]{ref-knuth1984a}
Knuth, D.E. 1984: {`Literate Programming'}. \emph{The Computer Journal}.
DOI:
\href{https://doi.org/10.1093/comjnl/27.2.97}{10.1093/comjnl/27.2.97}.

\hypertarget{ref-knuth1992a}{}

\bibitem[\citeproctext]{ref-knuth1992a}
Knuth, D.E. and Levy, S. n.d. {`The CWEB System of Structured
Documentation'}.
\url{https://www-cs-faculty.stanford.edu/~knuth/cweb.html}.

\hypertarget{ref-martin2017cleanArchitecture}{}

\bibitem[\citeproctext]{ref-martin2017cleanArchitecture}
Martin, R.C. 2017: \emph{Clean Architecture: A Craftsman's Guide to
Software Structure and Design}, 1st edn. Prentice Hall.

\end{CSLReferences}

\subsection*{B. Direct literate-programming tools and descendants}
\begin{CSLReferences}{1}{0}
\hypertarget{ref-gnuEmacsDownload}{}
\bibitem[\citeproctext]{ref-gnuEmacsDownload}
{`GNU Emacs Download'} n.d.
\url{https://www.gnu.org/software/emacs/download.html}.

\hypertarget{ref-orgmanualExtractingSourceCode}{}

\bibitem[\citeproctext]{ref-orgmanualExtractingSourceCode}
{`Org Manual: Extracting Source Code'} n.d.
\url{https://orgmode.org/manual/Extracting-Source-Code.html}.

\hypertarget{ref-ramsey1994a}{}

\bibitem[\citeproctext]{ref-ramsey1994a}
Ramsey, N. n.d. {`Noweb Repository'}.
\url{https://github.com/nrnrnr/noweb}.

\end{CSLReferences}

\subsection*{C. Executable documents / reproducible research / notebook lineage}
\begin{CSLReferences}{1}{0}
\hypertarget{ref-knitrDocs}{}
\bibitem[\citeproctext]{ref-knitrDocs}
{`Knitr'} n.d. \url{https://yihui.org/knitr/}.

\hypertarget{ref-nbconvertDocs}{}

\bibitem[\citeproctext]{ref-nbconvertDocs}
{`Nbconvert Documentation'} n.d.
\url{https://nbconvert.readthedocs.io/}.

\hypertarget{ref-pandocDocs}{}

\bibitem[\citeproctext]{ref-pandocDocs}
{`Pandoc'} n.d. \url{https://pandoc.org/}.

\hypertarget{ref-jupyterDocs}{}

\bibitem[\citeproctext]{ref-jupyterDocs}
{`Project Jupyter Documentation'} n.d. \url{https://docs.jupyter.org/}.

\hypertarget{ref-quartoDocs}{}

\bibitem[\citeproctext]{ref-quartoDocs}
{`Quarto Documentation'} n.d. \url{https://quarto.org/docs/}.

\hypertarget{ref-rmarkdownDocs}{}

\bibitem[\citeproctext]{ref-rmarkdownDocs}
{`R Markdown Documentation'} n.d. \url{https://rmarkdown.rstudio.com/}.

\end{CSLReferences}

\subsection*{D. LLM / AI-assisted programming}
\begin{CSLReferences}{1}{0}
\hypertarget{ref-vscodeCopilotCustomInstructions}{}
\bibitem[\citeproctext]{ref-vscodeCopilotCustomInstructions}
{`Customize GitHub Copilot in Visual Studio Code'} n.d.
\url{https://code.visualstudio.com/docs/copilot/copilot-customization}.

\hypertarget{ref-githubCopilotCustomization}{}

\bibitem[\citeproctext]{ref-githubCopilotCustomization}
{`Customizing GitHub Copilot'} n.d.
\url{https://docs.github.com/en/copilot/how-tos/custom-instructions}.

\end{CSLReferences}

\subsection*{E. Governance / provenance (contextual)}
\begin{CSLReferences}{1}{0}
\hypertarget{ref-c2paSpec}{}
\bibitem[\citeproctext]{ref-c2paSpec}
{`C2PA Specification'} n.d. \url{https://c2pa.org/specifications/}.

\hypertarget{ref-euCodeOfPracticeGpai}{}

\bibitem[\citeproctext]{ref-euCodeOfPracticeGpai}
{`General-Purpose AI Code of Practice'} n.d.
\url{https://digital-strategy.ec.europa.eu/en/policies/ai-code-practice}.

\hypertarget{ref-euAiAct}{}

\bibitem[\citeproctext]{ref-euAiAct}
\emph{Regulation (EU) 2024/1689 (Artificial Intelligence Act)} n.d.
\url{https://eur-lex.europa.eu/eli/reg/2024/1689/oj}.

\end{CSLReferences}

\end{document}
